%-----------------------------------------------
%  Engineer's Thesis
%  Copyright (c) 2026, Tomasz Nejman
%  Warsaw University of Technology, 2026
%-----------------------------------------------

\documentclass[
    bindingoffset=5mm,
    footnoteindent=3mm,
    hyphenation=true
]{src/wut-thesis}

\usepackage{xcolor}
\addbibresource{bibliografia.bib}

\newcommand{\refneeded}[1]{\textsuperscript{\textcolor{red}{#1}}}

\facultyeiti
\EngineerThesis
\langpol

\begin{document}

\instytut{Informatyki}
\kierunek{Informatyka}
\specjalnosc{Inżynieria Wytwarzania Oprogramowania}
\title{
    Projekt i implementacja komponentowego silnika \\
    gry strategicznej z turowym systemem rozgrywki
}
\engtitle{
    Design and implementation of a component-based \\
    strategy game engine with turn-based gameplay
}
\author{Tomasz Nejman}
\album{331411}
\promotor{mgr inż. Witold Wysota}
\date{2027}
\maketitle

\clearpage
\pagestyle{plain}
\begin{singlespace}
    \abstract Treść streszczenia po polsku (maksymalnie jedna strona).
    \keywords XXX, XXX, XXX
\end{singlespace}

\clearpage
\begin{singlespace}
    \secondabstract Abstract text in English (one page at most).
    \secondkeywords XXX, XXX, XXX
\end{singlespace}

\clearpage
\makeauthorship

\clearpage
\makelicense

\clearpage
\tableofcontents

\clearpage
\pagestyle{headings}

\section{Wstęp}

\subsection{Cel pracy inżynierskiej}

\subsection{Układ pracy}

\subsection{Słownik pojęć i oznaczeń}

\clearpage
\section{Wprowadzenie do architektury silników gier}

\subsection{Granica między silnikiem a grą}

Zrozumienie współczesnej architektury silników gier wymaga w pierwszej kolejności
zdefiniowania granicy między samym silnikiem a logiką gry. Jak zauważa Jason
Gregory w monografii \emph{Game Engine Architecture}~\cite{GEA}, w praktyce rozróżnienie, gdzie
kończy się silnik, a zaczyna gra bywa trudne, a granica pozostaje rozmyta. Do tego
silniki gier rzadko powstają w próżni. Przykładowo silnik \emph{Unreal Engine 1} (rok 1998)
tworzony był specyficznie dla flagowej gry Epic Games \emph{Unreal}.\refneeded{REF NEEDED}

\subsection{Od układów cyfrowych do monolitycznego kodu}

Historycznie, projektanci pierwszych gier nie rozróżniali pojęć “gra” i “silnik” w ogóle,
gdyż nie było możliwości myślenia na takim poziomie abstrakcji. Początkowo gry nie
były w ogóle “pisane” we współczesnym znaczeniu, z wykorzystywaniem języka
programowania. Analiza opisu patentowego \emph{Video Image Positioning Control System
    for Amusement Device}~\cite{PONG} wykazuje, iż maszyny odpowiedzialne za uruchomienie
gry \emph{PONG} (rok 1972) nie posiadały procesora, pamięci operacyjnej, a tym bardziej
kodu w dzisiejszym rozumieniu. Cała “logika” gry, tj. fizyka piłki, wyświetlanie punktów,
była permanentnie utrwalona na płytce TTL i jej układach cyfrowych. Mechanika gry
była de facto tożsama ze sprzętem komputerowym.

Sytuacja uległa zmianie wraz z upowszechnieniem się mikroprocesorów. Poziom
abstrakcji wzrósł ponad okablowanie do kodu maszynowego, jednak zasoby sprzętowe
były drastycznie ograniczone, co uwidacznia poniższa tabela:

\begin{table}[!h] \centering
    \caption{Ilość pamięci RAM wybranych urządzeń.}
    \begin{tabular}{| l | c | c |} \hline
        Urządzenie                     & Ilość pamięci RAM & Rok wydania \\ \hline\hline
        Oryginalny automat z PAC-MANem & 2 KB              & 1980        \\ \hline
        Commodore 64                   & 64 KB             & 1982        \\ \hline
        Atari 600XL                    & 64 KB             & 1983        \\ \hline
    \end{tabular}
\end{table}

W konsekwencji tych ograniczeń programy, a w szczególności gry, stanowiły
architektoniczny monolit. Przykładowo w grze \emph{PAC-MAN}~\cite{PACMAN} logika gry, AI
“duszków”, rysowanie grafiki i obsługa joysticka są ze sobą całkowicie wymieszane, ze
względu na wymóg intensywnej optymalizacji. Każda gra była tak naprawdę swoim
własnym, unikalnym silnikiem, którego nie dało się podzielić na komponenty, czy w
skuteczny sposób użyć ponownie.

\subsection{Wzrost mocy obliczeniowej i separacja silnika od gry}

Już w 1965 roku Gordon Moore~\cite{MOORE} zaobserwował wykładniczy wzrost liczby
tranzystorów w układach scalonych i przewidział, że utrzyma się on jeszcze przez
wiele lat. Ów wzrost spowodował monumentalne i stabilne zwiększanie się mocy
obliczeniowej komputerów/konsol i rozmiarów pamięci. Pozwoliło to na pojawienie się
gier, wpierw pseudo-3D\refneeded{REF P3D}, np. \emph{Wolfenstein 3D} (rok 1992) czy \emph{DOOM} (rok 1993), a
później true-3D\refneeded{REF T3D}, np. \emph{Quake} (rok 1996 - PC), \emph{Super Mario 64} (rok 1996 - Nintendo
64). Oprócz tego zwiększyły się możliwości nie tylko graficzne, ale także logiczne.
Gwałtowny wzrost liczby symulowanych obiektów oraz wprowadzenie przestrzeni
trójwymiarowej wywołały kryzys związany z zarządzaniem globalnym stanem gry. W
grach dwuwymiarowych ten stan mógł być kontrolowany przez relatywnie prostą
strukturę danych. W środowisku 3D, gdzie jednocześnie musiały współistnieć systemy
detekcji kolizji, kalkulacji fizyki oraz zaawansowanej sztucznej inteligencji, utrzymanie
spójności danych stało się wyzwaniem. Tradycyjne, monolityczne podejście, w którym
cały program miał bezpośredni dostęp do zmiennych globalnych, prowadziło do błędów
synchronizacji, trudnych do wykrycia efektów ubocznych oraz uniemożliwiało
skuteczne skalowanie kodu. To drastyczne skomplikowanie systemów zarządzania
stanem wymusiło ostateczną separację logiki od kodu bazowego. Ponownie
programiści musieli wejść na jeszcze wyższy poziom abstrakcji. Oprócz tego, jak
wspominają Bishop et al.~\cite{BISHOP} (1998) tworzenie tak skomplikowanych gier od zera
nie byłoby ani ekonomicznie opłacalne, ani konieczne. Definiują oni silnik gry jako zbiór
elementów gry, nie mających wpływu na samą zawartość gry, tj:
\begin{itemize}
    \item pobieranie informacji z urządzeń wejścia,
    \item symulowanie dynamiki obiektów,
    \item renderowanie obrazu,
    \item odtwarzanie dźwięku.
\end{itemize}

\subsection{Warstwy silnika i pętla gry}

\subsection{Zarządzanie stanem gry}

\subsection{Architektura komponentowa}

\subsection{Silniki gier strategicznych a uniwersalność kodu}

O ile powyższa taksonomia odpowiada klasycznym silnikom, szczególnie gier typu
FPS, TPS, MMO, MOBA, w przypadku gier strategicznych sytuacja wygląda nieco
inaczej. W przypadku gier RTS lub turowych priorytety przesuwają się z symulacji
ciągłej fizyki i predykcji stanów, w stronę deterministycznego zarządzania stanem wielu
obiektów, np. obsługi unikalnych wymagań strategii, takich jak zarządzanie armiami,
czy specyficzne algorytmy wyszukiwania ścieżek. Ponadto, wiele gier strategicznych
bazuje na planszach o dyskretnym podziale na pola, np. serie \emph{Civilization}
(sześciokątne pola), \emph{Heroes of Might and Magic} (kwadratowe i sześciokątne pola). Ten
gatunkowy rozstrzał bezpośrednio ilustruje problematykę uniwersalności kodu.
Problem ten definiuje Jason Gregory w monografii \emph{Game Engine Architecture}~\cite{GEA},
gdzie określa silnik gry jako oprogramowanie, które jest “możliwe do rozbudowania i
może być użyte jako podstawa dla wielu gier bez większych modyfikacji”.
Jednocześnie Gregory zaznacza, że w wielu przypadkach granica między silnikiem a
grą jest rozmyta. Podkreśla on, że przesunięcie tej granicy w jedną czy drugą stronę
ma bardzo jasne konsekwencje: silnik ogólnego przeznaczenia będzie bardziej
uniwersalny, ale mniej przystosowany do konkretnego typu gry, zaś silnik bardziej
nakierowany na dany podgatunek będzie znacznie trudniejszy lub wręcz niemożliwy do
użycia w przypadku ogólnym. Ta dwojaka natura silników powoduje, że projektowanie
architektury silnika to nieustanne zarządzanie kompromisem, z jednej strony
elastyczność z drugiej optymalizacja pod konkretny typ rozgrywki.

\clearpage
\section{Planowanie w grach turowych}

\subsection{Złożoność gier turowych}

Przede wszystkim warto wspomnieć, że w grach turowych nie sposób jest zbadać
wszystkie możliwe ścieżki gry. Na przykładzie klasycznej gry w szachy Claude E.
Shannon w pracy \emph{Programming a Computer for Playing Chess}~\cite{SHANNON} wylicza liczbę
możliwych rozgrywek na rząd wielkości $10^{120}$, z założeniem średnio 40 ruchów per
rozgrywka. Następnie z założeniem maszyny, która symuluje jedną rozgrywkę w czasie
jednej mikrosekundy oblicza średni czas na wybranie pierwszego optymalnego ruchu
na $10^{90}$ lat. Jest to liczba absolutnie przekraczająca możliwości obliczeniowe
współczesnych systemów.

Nowoczesne strategiczne gry turowe przewyższają szachy w możliwej liczbie decyzji,
ruchów i tur praktycznie zawsze. W złożonych grach typu 4X (eXplore, eXpand,
eXploit, eXterminate), takich jak \emph{Civilization VI}, standardowa rozgrywka wymaga
rozegrania setek tur (średnio ok 350~\cite{CIV}), a podczas każdej z nich gracz podejmuje
wiele decyzji (średnio 10-20), każda z nich mająca wiele różnych wyników. Wskutek
tego obliczenie optymalnego ruchu można uznać za niemożliwe w sensownym czasie.

\subsection{Minimax i przycinanie alfa-beta}

\subsection{Monte Carlo Tree Search i Statistical Forward Planning}

\subsection{Losowość i niepełna informacja}

\subsection{Paradygmat Forward Model}

Algorytmy sztucznej inteligencji bazują na rozgrywaniu tysięcy, jeśli nie milionów gier,
symulując podejmowanie decyzji i ewaluując nowy stan. W związku z tym, dostarczany
silnik musi spełniać paradygmat \emph{Forward Model}, tj. musi dostarczać deterministyczną
funkcję stan-przejście. Mając dany Stan (\emph{S}) i Akcję (\emph{A}) funkcja oblicza nowy stan
(\emph{S+1}). Ze względu na gigantyczną liczbę symulacji funkcja ta musi być możliwie
zoptymalizowana pod względem złożoności obliczeniowej i pamięciowej. Oprócz tego
wiele algorytmów drzewiastych (Minimax, MCTS) wymagają częstego duplikowania
obecnego stanu gry na poddrzewa. Jednakże kopiowanie stanu gry często bywa
bardzo kosztowne; przykładowo gra \emph{StarCraft II} przez wiele lat nie dostarczała API
typu \emph{Forward Model} właśnie ze względu na złożoność obliczeniową i pamięciową\refneeded{REF STARCRAFT}.
Warto dodać, że optymalny \emph{Forward Model} powinien projektować stan jako
dane ciągłe w pamięci, co pozwala na klonowanie stanu za pomocą prostych operacji
kopiowania bloków (cache hit) zamiast kosztownych operacji po obiektach.

\subsection{Strategie obsługi stanu: kopiowanie, cofanie, współdzielenie}

\clearpage
\section{Istniejące rozwiązania}

\subsection{Zakres przeglądu}

Mając na celu dostarczenie szybkie i wygodne FM API dla agentów AI oraz w związku
z naciskiem na gatunek TBS (Turn Based Strategy) gry \emph{Heroes of Might and Magic},
podczas rozpatrywania istniejących rozwiązań rynkowych, prac badawczych i
dedykowanych silników będę skupiał się nie na skomplikowanych potokach
renderowania, synchronizacji sieciowej czy algorytmach zbierania danych od
użytkownika, ale głównie na tym, w jaki sposób silniki te w deterministyczny sposób
zarządzają stanem gry. Co więcej, ze względu na charakter badawczy projektu mojego
silnika będę również skupiał się nie tylko na możliwości prowadzenia faktycznej
rozgrywki gracz-gracz, ale także na łatwości wykorzystania silnika przez agentów, na
cele treningu.

\subsection{Środowiska badawcze}

\subsubsection{Tribes}

Zagadnienie dostarczenia gry z API typu FM zostało zbadane i opisane przez Diego
Perez-Liebana et al. (2020) w pracy \emph{Tribes: A New Turn-Based Strategy Games for AI
    Research}~\cite{TRIBES}. Jest to prosta gra turowa typu 4X, bazująca na grach typu
\emph{Civilization}, \emph{The Battle of Polytopia}. Gra ta odbywa się na planszy 2D o rozmiarze N x
N. W każdej turze gracze mogą poruszać się po planszy, atakować, rozwijać drzewo
umiejętności, etc. Powodem stworzenia tej gry był właśnie narzut związany z
kopiowaniem stanu w klasycznych grach 4X (i klasycznych silnikach) na rzecz
trenowania agentów AI, a celem było stworzenie lekkiego mechanizmu replikowania
dla algorytmów typu SFP (ang. \emph{Statistical Forward Planning}). Oprócz tego gra dotyka
aspektu częściowej możliwości obserwowania stanu gry, tj. z perspektywy danego
gracza część planszy jest ukryta za tzw. “mgłą wojny”. Środowisko gry \emph{Tribes} z
powodzeniem realizuje interfejs \emph{Forward Model} udostępniając metody do kopiowania
stanu gry, wykonywania ruchów. Aby zapewnić czystszy interfejs, \emph{Tribes} wykorzystuje
wzorzec \emph{Command Pattern}, aby dokonać enkapsulacji i modularyzacji możliwych i
wykonywanych ruchów. Choć nie zaimplementowano tego w \emph{Tribes}, \emph{Command Pattern}
umożliwia również kolejkowanie ruchów lub implementację metody “undo”, jeżeli ruchy
są deterministyczne.

Perez-Liebana et al. (2020) skupia się głównie w \emph{Tribes} na kombinatorycznym
przyroście związanym z decyzjami globalnymi, tj. co wybudować, co zbadać, i
dalekosiężnym planowaniem (SFP) w warunkach niepełnej informacji.

\subsubsection{MicroRTS}

Problemu wykładniczego przyrostu poddrzew w algorytmach drzewiastych dotyka
również Santiago Ontañón w pracy \emph{The Combinatorial Multi-Armed Bandit Problem
    and Its Application to Real-Time Strategy Games}~\cite{MULTI_ARMED_BANDIT}. Praca ta
rozpatruje, jak sama nazwa wskazuje, problem \emph{Combinatorial Multi-Armed Bandit}. W
sytuacji gdy istnieje \emph{N} jednostek i każda z nich ma \emph{M} możliwych ruchów, to przestrzeń
akcji w jednym kroku wynosi $M^N$. Ontañón skupia się na próbie oszacowania
optymalnego ruchu poprzez rozważanie sposobów przeszukiwania przestrzeni i
balansu eksploracja-eksploracja na podstawie autorskiego silnika gry \emph{MicroRTS}.

Podobnie jak \emph{Tribes}, \emph{MicroRTS} dostarcza interfejs \emph{Forward Model}, czyli zapewnia
możliwość kopiowania stanu gry globalnym \emph{copy} i implementuje \emph{Command Pattern}. W
przeciwieństwie do \emph{Tribes}, \emph{MicroRTS}, jak sama nazwa wskazuje jest grą typu RTS
(Real Time Strategy), czyli czas w grze jest pseudo-ciągły (tick , a gracze podejmują
decyzje w dowolnym momencie, asynchronicznie.

O ile \emph{Tribes} optymalizuje “lekkość” samego stanu, by operacja \emph{copy} była jak najmniej
kosztowna pamięciowo i wykorzystuje płaskie struktury danych oraz prymitywy, o tyle
\emph{MicroRTS} często radzi sobie z problemem poprzez abstrakcję akcji, tj. skryptowanie
zachowań jednostek, aby odciążyć \emph{Forward Model}. W \emph{MicroRTS} nie planuje się kroków
co tick (znacznie wyższa granularność czasowa RTS), zamiast tego algorytm planuje
akcje długoterminowe.

\subsubsection{Ograniczenia środowisk opartych na Javie}

Środowiska badawcze, takie jak opisywane wcześniej \emph{Tribes} czy \emph{MicroRTS}, z
powodzeniem realizują interfejs \emph{Forward Model}, ich implementacja opiera się na
wysokopoziomowym, obiektowym języku programowania jakim jest Java. W Javie
wszystkie obiekty klas tworzone są na stercie, a zarządzanie pamięci operate jest o
garbage collector. W konsekwencji obydwa mechanizmy kopiowania stanu gry są
wysoce nieoptymalne, gdyż każde głębokie kopiowanie stanu gry skutkuje kaskadą
braków trafienia w pamięci podręcznej (cache miss)\refneeded{REF NEEDED}.

\subsubsection{Stratega}

\subsection{Rozwiązania dla Heroes of Might and Magic III}

\subsubsection{VCMI}

Oryginalna gra \emph{HoMM III} była napisana głównie w języku C++~\cite{ICCULUS}, ale kod do tej
pory jest prywatną własnością (obecnie) firmy Ubisoft. \emph{VCMI} to projekt open-source,
mający na celu zbudowanie gry \emph{Heroes of Might and Magic III} od zera w języku C++~\cite{VCMI}.
Silnik \emph{VCMI} jest budowany z myślą o replikacji gry 1 do 1, z podejściem OOP.
\emph{VCMI} nie jest jednolitym programem, dziel się na proces serwera, który autoryzuje stan
gry, oraz klientów odpowiedzialnych za renderowanie i obsługę gracza. Taka separacja
jest dobrym rozwiązaniem do rozgrywek sieciowych, ale dla agenta AI oznaczałby duży
narzut czasowy związany z tym, że komunikacja między botem a serwerem często
musi przechodzić przez warstwę sieciową. Część serwerowa \emph{VCMI} jest ściśle
ustrukturyzowana i zorientowana na dostęp do obiektów przez czyste interfejsy, jednak
to wprowadza problem nieciągłości danych w pamięci. Twórcy \emph{VCMI} nie skupiali się na
tym, gdyż głównym celem było utworzenie działającej i utrzymywalnej wersji gry
\emph{Heroes III}, a nie środowiska dla agentów AI. W konsekwencji tego \emph{VCMI} nie
wprowadza interfejsu \emph{Forward Model}.

\subsubsection{vcmi-gym i MMAI}

\subsection{Silniki komercyjne}

\subsection{Porównanie rozwiązań i wskazanie luki}

\clearpage
\section{Analiza wymagań}

\subsection{Odwzorowane i pominięte mechaniki gry}

\subsection{Wymagania funkcjonalne}

\subsection{Wymagania niefunkcjonalne}

\subsection{Wybór technologii}

\clearpage
\section{Projekt silnika}

\subsection{Architektura komponentowa silnika}

\subsection{Oddzielenie logiki gry od renderowania i obsługi wejścia}

\subsection{Maszyna stanów gry}

\subsection{Polecenia jako dane}

\subsection{Interfejs Forward Model}

\subsection{Projekt stanu bitwy}

\subsection{Interfejs agentów}

\clearpage
\section{Implementacja}

\subsection{Środowisko i proces budowania}

\subsection{Stan bitwy i jego kopiowanie}

\subsection{Generowanie legalnych poleceń i przejścia między stanami}

\subsection{Integracja agentów}

\subsection{Warstwa prezentacji}

\subsection{Testy i weryfikacja}

\clearpage
\section{Podsumowanie}

\subsection{Wnioski}

\subsection{Kierunki rozwoju}

\clearpage
% \nocite{*}
\printbibliography

\clearpage
\listoffigurestoc

\clearpage
\listoftablestoc

\clearpage
\listofappendicestoc

\clearpage
\appendix{Instrukcja uruchomienia}

\end{document}
